% Source: https://artofproblemsolving.com/community/c3h2790795_2019_mbmt_guts_round_d115_l19_montgomery_blair_math_tournament
% Stable identity: aops-topic-2790795-q14
Set $ A=(0,0,0),B=(1,1,1)$. A point on a face not incident to $ A $ has one coordinate 1, so Euclidean distance at least 1, with equality only at a neighbouring vertex. Thus every permissible $ X $ lies on a face incident to $ A $. By symmetry write $ X=(x,y,0)$ with $ x,y\ge0 $ and $ x^2+y^2=1 $; the straight face segment from $ A $ gives this exact condition. Unfolding the bottom face with either adjacent side gives valid paths from $ B $ of lengths $\sqrt{(2-x)^2+(1-y)^2}$ and $\sqrt{(1-x)^2+(2-y)^2}$. Assume $ x\ge y $; the first squared length is $ 6-4x-2y $. Since $(2x+y)^2=4-3y^2+4xy\ge4 $, we have $ 2x+y\ge2 $, so that path has length at most $\sqrt2 $. The straight unfolded segments cross their shared edges within the unit edge because their horizontal or vertical coordinates interpolate between 0 and 1. Equality is attained at $ X=(1,0,0)$: it and $ B $ lie on a common face at Euclidean distance $\sqrt2 $, which no surface path can improve. Hence the maximum is $\boxed{\sqrt2}$.
\begin{asy}
usepackage("amsmath");
size(7cm);
draw((0,0)--(2,0)--(2,1)--(0,1)--cycle);draw((1,0)--(1,1),dashed);pair A=(0,0),X=(0.8,0.6),B=(2,1);draw(A--X,blue);draw(X--B,red);label("$A$",A,SW);label("$X$",X,S);label("unfolded $B$",B,NE);label("bottom face",(0.5,1),N);label("side face",(1.5,1),N);
\end{asy}
